\chapter{Standard og strid}

\begin{motto}
Skulle forma kome frå typen eller frå kunstnaren?\\ Spørsmålet vart aldri svara på. Det vart røysta over, og deretter gløymt.
\end{motto}

Der Arts and Crafts vende ryggen til maskinen, vende Deutscher Werkbund seg mot han. Stifta i München i 1907 av ei blanding av arkitektar, kunstnarar, handverkarar og — det er det nye — industrimenn og politikarar, var Werkbund det fyrste store forsøket på å gje formgjevinga eit grunnlag som rekna med fabrikken i staden for mot han. Ambisjonen var nasjonaløkonomisk så vel som estetisk: tysk industri skulle vinne verdsmarknaden ikkje ved å vere billegast, men ved å vere best forma, og for å bli best forma måtte han alliere seg med kunsten. Her finn me, for fyrste gong, draumen om at form kunne bli noko meir enn smak — at han kunne bli ein standard, noko etterprøvbart, noko ein kunne einast om og produsere i tusental. Og her finn me, alt i opningsåra, oppdaginga av at faget ikkje hadde nokon teori til å avgjere det mest grunnleggjande spørsmålet sitt. Striden kom, og striden vart ikkje avgjord. Han vart røysta over.

\section{Muthesius og det engelske førelegget}

Mannen som meir enn nokon forma Werkbund sin tidlege ambisjon, kom til saka via England. Hermann Muthesius hadde vore knytt til den tyske ambassaden i London gjennom store delar av nittiåra og dei fyrste åra etter hundreårsskiftet, med eit oppdrag som var halvt diplomatisk og halvt etterretningsfagleg: å studere engelsk bustadbygging og bringe lærdomen heim. Resultatet vart trebandsverket \textit{Das englische Haus}, der \textcite{muthesius1904} legg fram det engelske huset — og særleg arven etter Arts and Crafts — for eit tysk publikum med ein grundigheit ingen engelskmann hadde mønstra om sitt eige. Det er ein paradoksal bru. Den tyske moderniteten lærte mykje av sin sakleg-bustadlege grunnhaldning frå Morris og krinsen hans, formidla gjennom Muthesius. Men sjå kva Muthesius gjorde med arven: han skilde han frå handverksmoralen.

Det Muthesius beundra ved det engelske huset, var saklegheita, det føremålstenlege, fråvêret av påklistra historisme — den reine, ærlege bruksforma. Det han ikkje tok med seg, var Morris si overtyding om at denne forma berre kunne kome frå fritt handverk. Tvert om: Muthesius meinte at den saklege forma kunne, og burde, realiserast industrielt, i serie, etter standardiserte typar. Han trekte med andre ord den estetikken Arts and Crafts hadde grunna i ein arbeidsmoral, og kopla han frå den moralen — og kopla han i staden til fabrikken og marknaden. Dette er ein operasjon av stor rekkjevidd, og den var stilt mogleg nett fordi grunnlaget i Arts and Crafts var laust. Hadde rørsla hatt ein stram teori som batt forma uløyseleg til produksjonsmåten, kunne ein ikkje berre ha lyfta estetikken ut og sett han i fabrikken. Men sidan bandet mellom form og arbeid var ein moral og ikkje ein mekanisme, lét estetikken seg flytte fritt. Muthesius arva stilen utan moralen, akkurat slik me såg han var det einaste overførbare.

\section{Typisierung: striden i 1914}

Spenninga som låg under heile Werkbund-prosjektet, brast opent på årsmøtet i Köln sommaren 1914, vekene før krigen. Muthesius la fram ti teser, og kjernen i dei var ordet \emph{Typisierung}: typedanning, standardisering. Form måtte, hevda han, utviklast mot allmenngyldige typar — kanoniske, etterprøvde løysingar som kunne produserast i serie, eksporterast og kjennast att som tysk kvalitet. Berre gjennom typen kunne formgjevinga bli noko meir enn den einskilde kunstnaren sitt innfall; berre gjennom typen kunne ho bli ein delt standard, ein felleseige, noko ein industri og ein nasjon kunne byggje på. Det er, om ein ser nøye etter, eit forsøk på akkurat det denne boka spør etter: eit overførbart, ikkje-personleg grunnlag for forma.

Mot Muthesius reiste Henry van de Velde seg, og med han ein flokk av kunstnarane. Van de Velde sine motteser forsvara det motsette: kunstnaren sin individualitet, retten og plikta til alltid å søkje den nye forma, motstanden mot all standardisering som ein trussel mot det skapande. Ingen sann kunstnar, heldt han, ville nokon gong underordne seg ein kanon av typar; det skapande er per natur det som bryt typen. Striden var ekte og hard, og ho delte rørsla på midten. Den unge Walter Gropius stod på van de Velde si side — eit faktum me skal hugse når neste kapittel fører oss til Bauhaus, for det viser at sjølv grunnleggjaren av den mest «industrielle» av skular hadde sine røter i forsvaret for kunstnaren mot standarden. Striden gjekk til røysting, og enda i praksis uavgjort, overskugga få veker seinare av at verda gjekk i krig.

Sjå no nøye på kva dette var, for det er kjernen i kapitlet. Typisierung-striden var ikkje ein fagleg usemje som kunne avgjerast ved fagleg argumentasjon, slik to fysikarar avgjer ein usemje ved eit eksperiment. Det fanst ingen prøve som kunne ha vist at Muthesius hadde rett og van de Velde feil, eller omvendt. Begge posisjonane var koherente; begge hadde gode grunnar; og det fanst ikkje noko felles grunnlag dei begge kjende seg bundne av, som kunne ha avgjort saka. Difor måtte ho gå til røysting — og ei røysting er nettopp det ein gjer når ein ikkje har ein teori som avgjer. Spørsmålet «skal forma kome frå typen eller frå kunstnaren» er eit verdispørsmål forkledd som eit fagleg, og faget hadde — og fekk aldri — noko som kunne gjere det om til eit fagleg som let seg avgjere. Det \textcite{schwartz1996werkbund} viser med stor presisjon, er at heile Werkbund-debatten frå fyrst til sist dreidde seg om tilhøvet mellom form og massekultur, mellom merkevare og standard, mellom kunst og marknad — og at desse spenningane ikkje vart løyste, men levde vidare som ein varig uro i sjølve prosjektet om moderne form.

\section{Loos og det andre angrepet på ornamentet}

Parallelt med Werkbund, og i ein heilt annan tone, kom det andre store framstøytet mot den gamle forma — denne gongen frå Wien, frå arkitekten Adolf Loos. I essayet som er kjent som «Ornament og forbrytelse», halde som førelesing kring 1910 og trykt seinare, gjer \textcite{loos1908ornament} eit argument som ved fyrste augnekast liknar Werkbund sitt, men som kviler på eit heilt anna og langt meir tvilsamt grunnlag. Loos hevdar at ornament er eit teikn på primitivitet, og at det moderne, siviliserte mennesket har vakse forbi det. Tatoveringa høyrer papuanaren og forbrytaren til; det same gjer den utsmykka flata. Framsteget i kulturen, seier Loos, er sjølve fjerninga av ornament frå bruksgjenstanden. Den glatte, reine flata er ikkje berre vakrare; ho er meir moralsk, meir økonomisk — ornamentet er bortkasta arbeid og bortkasta kapital — og framfor alt meir \emph{moderne}.

Argumentet er retorisk strålande og fagleg verdilaust, og det er verdt å seie kvifor, fordi det viser noko om heile epoken. Loos grunngjev eit estetisk val — det reine over det pynta — med ein teori om kulturell utvikling som er rein spekulasjon, og som i ettertid er direkte ubehageleg: ein rangstige av folkeslag der europearen står på toppen og dei tatoverte nederst. Det finst ikkje noko prov for at ornament korrelerer med primitivitet; påstanden er eit fordom kledd som ein lov. Men legg merke til strukturen, for han er den same som me har sett før og skal sjå att: ein smak — Loos sin smak for den glatte flata — vert presentert ikkje som ein smak, men som ein naudsynt konsekvens av noko utanfor smaken, her av sjølve sivilisasjonen sin framgang. «Eg likar ikkje ornament» vert til «ornament er forbrytelse». Den private preferansen får lånt autoritet frå ein pseudoteori, og blir dermed unndregen diskusjon. Du kan ikkje vere usamd i Loos sin smak utan å framstå som tilbakeståande.

\textcite{banham1960} gjev oss reiskapen til å sjå dette klart. I si analyse av teori og form i den fyrste maskinalderen viser Banham korleis modernismens pionerar gong på gong grunngav formval med argument som hevda å vere objektive — funksjonelle, tekniske, historiske, evolusjonære — men som ved nærare ettersyn var estetiske val utstyrte med ein teoretisk overbygnad som ikkje bar. Loos er eit reindyrka tilfelle. Banham viser at den «maskinestetikken» modernismen meinte å avlese frå sjølve maskinen og det moderne livet, i røynda var ein stilpreferanse som var vald og deretter rasjonalisert; pionerane såg ikkje forma i maskinen, dei tok med seg ein smak og las han inn. Dette er det same diagnostiske grepet som låg under formelkritikken: forma følgjer ikkje funksjonen, men nokon sin smak som kallar seg funksjon — eller her, som kallar seg sivilisasjon.

\section{Behrens og AEG: standarden som vart ein merkevare}

Det finst eitt tilfelle frå Werkbund-krinsen der noko som likna ein standard faktisk vart realisert, og det er verdt å granske nøye, for det syner kva slag «standard» faget eigentleg var i stand til å produsere. I 1907 — same år som Werkbund vart stifta — tilsette den tyske elektrokonsernen AEG arkitekten Peter Behrens som kunstnarleg rådgjevar med eit oppdrag av ein type som ikkje hadde funnest før: han skulle gje heile konsernet eitt samla formuttrykk. Og det gjorde han. Behrens teikna AEG sine fabrikkbygningar, fremst turbinhallen i Berlin frå 1909; han teikna produkta, frå dei berømte elektriske vasskjelane til vifter og lampar; han teikna logoen, brevarka, katalogane, butikkane, plakatane. For fyrste gong i historia fekk eit selskap ein gjennomført visuell identitet frå bygning til brevhovud, utforma av éin styrande hand. Behrens vert difor ofte kalla den fyrste industridesignaren, og AEG-arbeidet det fyrste tilfellet av det me i dag kallar ein heilskapleg merkevareidentitet. Her, skulle ein tru, ser me Werkbund sin draum om standardisert kvalitetsform realisert i praksis, ikkje berre debattert i Köln.

Sjå likevel kva slag samanheng dette var, for det er avgjerande. AEG-forma hang saman — ho var verkeleg eit system, gjenkjenneleg på tvers av turbinhall og tekjel — men ho hang ikkje saman fordi det fanst ein \emph{teori} som sa korleis ein AEG-gjenstand skulle sjå ut. Ho hang saman fordi det fanst éin mann, Behrens, med eit auge og ein autoritet, som teikna eller godkjende alt. Samanhengen budde i personen, ikkje i eit fundament. Då Behrens gjekk, eller då ein annan konsern ville ha det same, fanst det ingen overførbar lære å sende med; ein måtte tilsetje ein ny Behrens, eit nytt styrande auge. Og legg merke til kva som batt forma saman på papiret: ikkje eit kvalitetskriterium nokon kunne etterprøve, men eit eigarskap — AEG sitt namn, AEG si merkevare. Det som gav forma einskap, var ikkje at ho var \emph{rett}, men at ho var \emph{AEG si}. Reyner Banham peikar på at Behrens nett difor står som ein nøkkelfigur: i han ser ein at den moderne, sakleg-industrielle forma frå starten var bunden til selskapsidentitet og marknad, ikkje til nokon utleidd teori om kva industriell form skulle vere \autocite{banham1960}.

\begin{kasus}{AEG 1907: den fyrste merkevara, ikkje den fyrste standarden}
Peter Behrens gav AEG alt det Muthesius drøymde om: ei samla, sakleg, attkjenneleg kvalitetsform, produsert i serie, eksportabel, knytt til tysk industri. Og likevel er AEG-tilfellet ikkje eit prov på at Werkbund fann standarden — det er det skarpaste provet på at det ikkje gjorde det. For det som heldt AEG-forma saman, var ikkje ein delt teori om industriell form som ein annan bedrift kunne ha teke i bruk; det var Behrens sitt eige auge og AEG sitt eige namn. Forma var ein \emph{merkevare}, ikkje ein \emph{standard}. Skilnaden er heile kapitlet: ein standard er eit kollektivt minne som høyrer faget til og kan brukast av alle, slik ein gjengestorleik kan; ein merkevare er ein eigedom som høyrer eitt selskap til og bind berre det. Werkbund kunne produsere det andre og aldri det fyrste, fordi det fyrste krev eit kriterium og eit språk faget ikkje hadde, medan det andre berre krev ein eigar og eit auge. Då nokon endeleg klarte å gje den industrielle forma einskap, viste det seg at einskapen heitte AEG — altså eit firmanamn, ikkje eit prinsipp.
\end{kasus}

Dette skuggar langt framover, heilt inn i vår eiga tid, og det er verdt å seie det med ein gong. Når designfaget i dag skal peike på sine fremste prestasjonar i å gje form einskap og konsistens, peikar det påfallande ofte på \emph{merkevarer} — på Braun under Rams, på Apple, på den gjennomførte identiteten til eitt selskap under ei styrande hand — og påfallande sjeldan på delte, faglege standardar som høyrer faget til og ikkje til ein eigar. Det er ikkje eit samantreff. Det er AEG-mønsteret som gjentek seg: den einaste forma faget har lært å gjere verkeleg konsistent, er den forma som ein eigar betalar for og ein autoritet held saman. Behrens er det fyrste tilfellet, og han set malen. Werkbund ville gje forma eit fundament og gav henne i staden ei merkevare — og forveksla, slik faget skulle halde fram med å forveksle, den eine med den andre.

\section{Standarden som mangla eit språk}

Det er verdt å dvele ved kva Werkbund eigentleg prøvde å byggje, fordi ambisjonen er den mest lovande me har møtt så langt, og fallet difor det mest opplysande. Ein standard er, rett forstått, eit kollektivt minne. Når ein bransje einar seg om ein type — ein gjengestorleik, eit papirformat, ein spenning i nettet — har han lagra ei løysing utanfor det einskilde hovudet, slik at neste generasjon ikkje treng å finne henne på nytt. Det er denne forma for kumulativ, ikkje-personleg kunnskap Muthesius drøymde om for forma: ein kanon av typeløysingar som var prøvde, godtekne og felles, slik at formgjevinga kunne byggje vidare i staden for å starte forfrå med kvar kunstnar. Hadde dette late seg gjere, ville designfaget fått sitt fyrste fundament i ordets eigentlege meining — eit minne lagra i typar, ikkje i meistrar.

Men ein standard krev to ting Werkbund ikkje hadde. Han krev eit kriterium for kva som tel som ei god løysing — ein måte å avgjere at denne typen er betre enn den, som ikkje berre er smaken til den som foreslår. Og han krev eit språk presist nok til at to fagfolk meiner det same med dei same orda, slik at dei kan einast om kva dei einar seg om. Industrien hadde begge delar for det tekniske: ein million skruar med same gjenge passar fordi «gjenge» tyder noko eksakt og «passar» kan målast. Werkbund hadde ingen av delane for det estetiske. Det fanst inga måling som kunne skilje den gode forma frå den dårlege, og det fanst inga fagspråk der orda — saklegheit, kvalitet, eigenart — tydde det same for Muthesius som for van de Velde. Difor kunne standarden lagast for det tekniske og aldri for det forma faget kalla sitt eige. Det Werkbund prøvde, var å standardisere noko det ikkje hadde eit språk til å spesifisere. Striden i Köln var det augneblinket dette vart synleg: ikkje ein usemje innanfor eit felles språk, men oppdaginga av at noko felles språk ikkje fanst.

Dette fråvêret av eit delt fagspråk er den andre store mangelen denne boka sporar, etter fråvêret av eit etisk organ, og Werkbund er staden der han fyrst trer fram i full breidd. Eit modent fag kjennest att på at orda det dømer med er skarpe nok til å binde. Werkbund ville byggje ein standard og oppdaga at det ikkje hadde orda til å skrive han. Det fall difor tilbake på det same faget alltid fell tilbake på når språket sviktar: autoriteten til personar. Striden vart ikkje avgjord av eit kriterium, men forsøkt avgjord av dei mest vektige stemmene i rommet, og då dei var usamde, av ei røysting. Eit fag som må røyste over forma si, har innrømt at det ikkje kan skrive henne ned.

\section{Eit verdival, ikkje ein avgjord teori}

Lat oss samle det Werkbund og Loos lærer oss om grunnlaget. Werkbund gjorde det store, ærlege forsøket på å flytte forma frå smak til standard — frå det den einskilde kunstnaren føretrekte, til ein delt, etterprøvd type som ein industri kunne byggje på. Det var rett ambisjon. Hadde det lukkast, ville faget hatt noko av det denne boka etterlyser: eit ikkje-personleg grunnlag, overførbart, kumulativt. Men i augneblinken då ambisjonen møtte motstand — då van de Velde reiste seg for kunstnaren mot typen — viste det seg at faget ikkje hadde ressursane til å avgjere saka fagleg. Det fanst ingen teori om forma som kunne seie kven som hadde rett. Det fanst berre to verdihaldningar, begge respektable, og ein generalforsamling som kunne røyste. Striden vart ikkje løyst; han vart utsett, og deretter overskugga av historia.

Dette er det avgjerande, og det er lett å oversjå fordi det ser ut som ein vanleg fagleg debatt. Men ein fagleg debatt i eit modent fag har ein avgjeringsprosedyre: eit eksperiment, eit prov, eit kriterium begge partar på førehand godtek som bindande. Typisierung-striden hadde ingen. Spørsmålet om standard mot individ var, og er framleis, eit val mellom verdiar — mellom det allmenne og det eineståande, det reproduserbare og det unike — og ikkje eit spørsmål med eit fasitsvar faget kunne nærme seg over tid. At striden enda i ei røysting og ikkje i ein konklusjon, er ikkje ein tilfeldig historisk detalj; det er ei måling av kva grunnlag faget hadde. Ein røystar om verdiar. Ein måler ikkje. Og fordi spørsmålet aldri vart avgjort, men berre lagt bort, kom det att — i Bauhaus, i Ulm, i kvar seinare strid mellom det kunstnarlege og det systematiske i formgjevinga.

At dette ikkje berre var ein abstrakt teoretisk svikt, men fekk verkelege historiske konsekvensar, kan ein sjå i kva som seinare vart gjort i namnet til den same forma. \textcite{betts2004} sporar korleis den saklege, standardiserte tyske bruksforma — arven frå Werkbund — vart boren vidare gjennom heilt ulike politiske regime og gjeve heilt ulike politiske meiningar, frå Weimar gjennom det tredje riket og inn i Vest-Tyskland sin etterkrigsidentitet. Den same reine forma kunne tene demokratisk fornying og autoritær orden; ho hadde ingen iboande politikk, fordi ho aldri hadde fått ein grunn som batt henne til noko. Forma var ledig til å tene den som tok henne i bruk. Eit fag som hadde grunngjeve forma si, ville hatt eit svar på kva ho var god for og kva ho ikkje skulle brukast til. Werkbund hadde gjeve forma ein prestisje og ein standard, men ingen samvit til å styre han — same fråvêret me såg hjå Arts and Crafts, no flytt inn i fabrikken.

\vignett

Werkbund kom nærare ein standard enn nokon før, og oppdaga ved den fyrste prøven at det ikkje hadde ein teori til å forsvare standarden mot kunstnaren. Striden vart utsett, ikkje avgjord. Men nokre få år etter krigen skulle ein ny skule i Weimar ta opp att akkurat denne uavgjorde striden — mellom kunst, handverk og industri — og denne gongen ikkje berre røyste over han, men byggje ein heil pedagogikk for å løyse han. Det vart det næraste designfaget nokon gong kom ein vitskap om form.

\laeringsmal{\begin{itemize}
\item Forklare kva \emph{Typisierung} var, og kvifor striden mellom Muthesius og van de Velde i 1914 ikkje kunne avgjerast fagleg, men måtte gå til røysting.
\item Skilje mellom ein \emph{standard} (kollektivt minne som høyrer faget til) og ein \emph{merkevare} (eigedom som bind eitt selskap), og bruke skiljet på Behrens og AEG.
\item Analysere Loos sitt ornamentargument som ein smak utstyrt med ein pseudoteori om kulturell utvikling.
\item Gjere greie for kvifor mangelen på eit delt fagspråk er den andre store mangelen boka sporar, etter mangelen på eit etisk organ.
\end{itemize}}

\ovingar{\begin{enumerate}
\item «Eit fag som må røyste over forma si, har innrømt at det ikkje kan skrive henne ned.» Finn ein samtidig situasjon der designarar «røystar» (panel, jury, avstemming) i staden for å måle. Kva fortel det om grunnlaget?
\item Behrens gav AEG ein merkevare, ikkje ein standard. Vel ein merkevare du kjenner godt og avgjer kva som held forma hennar saman: eit etterprøvbart prinsipp, eller ein eigar og eit auge? Kunne ein annan bruke same «standard»?
\item Industrien einar seg lett om tekniske standardar (gjenge, papirformat, spenning) men ikkje om estetiske. Diskuter kvifor. Er det ein språkmangel, ein verdimangel, eller begge?
\item Betts viser at den same saklege forma tente både demokrati og autoritært styre. Drøft: kan ei form ha ein iboande politikk, eller er ho alltid ledig til å tene den som tek henne i bruk?
\end{enumerate}}

\vidarelesing{\fullcite{schwartz1996werkbund}; \fullcite{muthesius1904}; \fullcite{loos1908ornament}; \fullcite{banham1960}; \fullcite{betts2004}; \fullcite{pevsner1936}.}
